% Source: rnn_gradient.tex
\begin{notechapter}
\chapter{Gradient Computation in RNNs}\label{ch:rnn-gradient}

\section{Introduction}
This document provides a detailed explanation of gradient computation in 
Recurrent Neural Networks (RNNs) using Backpropagation Through Time 
(BPTT). We start by defining the equations of the RNN, including the roles 
of the weight matrices and biases. We then split the general case into 
computations for the final step and computations for non-final steps, and 
illustrate the concepts with a concrete example of an RNN with \(\tau = 3\) 
time steps, fully developing all equations. We also include a detailed 
derivation of the gradient of the Negative Log-Likelihood Loss with respect 
to the output \(\mathbf{o}^{(t)}\), including all computational steps for 
the derivative of the softmax function.

\section{RNN Equations}
An RNN processes a sequence of inputs \(\mathbf{x}^{(1)}, \mathbf{x}^{(2)}, \ldots, \mathbf{x}^{(\tau)}\) 
and produces a sequence of hidden states \(\mathbf{h}^{(1)}, \mathbf{h}^{(2)}, \ldots, \mathbf{h}^{(\tau)}\) 
and outputs \(\mathbf{o}^{(1)}, \mathbf{o}^{(2)}, \ldots, \mathbf{o}^{(\tau)}\). The equations 
governing the RNN are as follows:

\subsection{Hidden State Update}
The hidden state \(\mathbf{h}^{(t)}\) at time step \(t\) is computed as:
\[
\mathbf{h}^{(t)} = \sigma(\mathbf{W}_{hh} \mathbf{h}^{(t-1)} + \mathbf{W}_{xh} \mathbf{x}^{(t)} + \mathbf{b}_h)
\]
where:
\begin{itemize}
    \item \(\mathbf{W}_{hh}\) is the hidden-to-hidden weight matrix.
    \item \(\mathbf{W}_{xh}\) is the input-to-hidden weight matrix.
    \item \(\mathbf{b}_h\) is the hidden state bias vector.
    \item \(\sigma\) is the activation function (e.g., \(\tanh\)).
\end{itemize}

\subsection{Output Computation}
The output \(\mathbf{o}^{(t)}\) at time step \(t\) is computed as:
\[
\mathbf{o}^{(t)} = \mathbf{W}_{ho} \mathbf{h}^{(t)} + \mathbf{b}_o
\]
where:
\begin{itemize}
    \item \(\mathbf{W}_{ho}\) is the hidden-to-output weight matrix.
    \item \(\mathbf{b}_o\) is the output bias vector.
\end{itemize}

\subsection{Loss Function}
The loss \(L\) is the sum of the losses at each time step:
\[
L = \sum_{t=1}^{\tau} L^{(t)}
\]
For classification tasks, the loss at each time step \(L^{(t)}\) is often the 
negative log-likelihood (cross-entropy) loss:
\[
L^{(t)} = -\sum_{i=1}^C y_i^{(t)} \log(\hat{y}_i^{(t)})
\]
where:
\begin{itemize}
    \item \(C\) is the number of classes.
    \item \(y_i^{(t)}\) is the true label (one-hot encoded).
    \item \(\hat{y}_i^{(t)}\) is the predicted probability for class \(i\).
\end{itemize}

\section{General Case: Gradient Computation in RNNs}
In an RNN, the hidden state \(\mathbf{h}^{(t)}\) at time step \(t\) 
influences the loss \(L\) in two ways:
\begin{enumerate}
    \item \textbf{Directly}: Through the current output \(\mathbf{o}^{(t)}\), 
          which contributes to the current loss \(L^{(t)}\).
    \item \textbf{Indirectly}: Through the next hidden state 
          \(\mathbf{h}^{(t+1)}\), which affects all future losses 
          \(L^{(t+1)}, L^{(t+2)}, \ldots\).
\end{enumerate}

The gradient computation differs for the final time step and non-final 
time steps. We address each case separately.

\subsection{Final Time Step (\(t = \tau\))}
At the final time step \(\tau\), the hidden state \(\mathbf{h}^{(\tau)}\) 
only affects the current loss \(L^{(\tau)}\). There is no future hidden 
state to consider. Thus, the gradient of the loss \(L\) with respect to 
\(\mathbf{h}^{(\tau)}\) is:
\[
\nabla_{\mathbf{h}^{(\tau)}} L = 
\left( \frac{\partial \mathbf{o}^{(\tau)}}{\partial \mathbf{h}^{(\tau)}} \right)^\top 
\left( \nabla_{\mathbf{o}^{(\tau)}} L \right)
\]
Here:
\begin{itemize}
    \item \(\frac{\partial \mathbf{o}^{(\tau)}}{\partial \mathbf{h}^{(\tau)}}\) 
          is the Jacobian matrix of the output with respect to the hidden 
          state. For a linear output layer, this is \(\mathbf{W}_{ho}\), so its transpose \(\mathbf{W}_{ho}^\top\) appears in the formula.
    \item \(\nabla_{\mathbf{o}^{(\tau)}} L\) is the gradient of the loss 
          with respect to the output \(\mathbf{o}^{(\tau)}\). For negative 
          log-likelihood loss, this is \(\hat{\mathbf{y}}^{(\tau)} - \mathbf{y}^{(\tau)}\).
\end{itemize}

\subsection{Non-Final Time Steps (\(t < \tau\))}
For non-final time steps (\(t < \tau\)), the hidden state \(\mathbf{h}^{(t)}\) 
affects both the current loss \(L^{(t)}\) and the next hidden state 
\(\mathbf{h}^{(t+1)}\). Thus, the gradient of the loss \(L\) with respect to 
\(\mathbf{h}^{(t)}\) is:
\[
\nabla_{\mathbf{h}^{(t)}} L = 
\left( \frac{\partial \mathbf{h}^{(t+1)}}{\partial \mathbf{h}^{(t)}} \right)^\top 
\left( \nabla_{\mathbf{h}^{(t+1)}} L \right) + 
\left( \frac{\partial \mathbf{o}^{(t)}}{\partial \mathbf{h}^{(t)}} \right)^\top 
\left( \nabla_{\mathbf{o}^{(t)}} L \right)
\]
Here:
\begin{itemize}
    \item \(\frac{\partial \mathbf{h}^{(t+1)}}{\partial \mathbf{h}^{(t)}}\) 
          is the Jacobian matrix of the hidden state transition. It describes 
          how \(\mathbf{h}^{(t+1)}\) changes with respect to \(\mathbf{h}^{(t)}\).
    \item \(\nabla_{\mathbf{h}^{(t+1)}} L\) is the gradient of the loss 
          with respect to \(\mathbf{h}^{(t+1)}\). It is computed recursively 
          during backpropagation.
    \item \(\frac{\partial \mathbf{o}^{(t)}}{\partial \mathbf{h}^{(t)}}\) 
          is the Jacobian matrix of the output with respect to the hidden 
          state. For a linear output layer, this is \(\mathbf{W}_{ho}\), so its transpose \(\mathbf{W}_{ho}^\top\) appears in the formula.
    \item \(\nabla_{\mathbf{o}^{(t)}} L\) is the gradient of the loss 
          with respect to the output \(\mathbf{o}^{(t)}\). For negative 
          log-likelihood loss, this is \(\hat{\mathbf{y}}^{(t)} - \mathbf{y}^{(t)}\).
\end{itemize}

\section{Negative Log-Likelihood Loss}
The negative log-likelihood loss is commonly used for classification tasks. 
It measures the discrepancy between the predicted output \(\hat{\mathbf{y}}^{(t)}\) 
and the true target \(\mathbf{y}^{(t)}\).

\subsection{Definition}
The negative log-likelihood loss for a single example is given by:
\[
L^{(t)} = -\sum_{i=1}^C y_i^{(t)} \log(\hat{y}_i^{(t)})
\]
where:
\begin{itemize}
    \item \( C \) is the number of classes.
    \item \( y_i^{(t)} \) is the true label (one-hot encoded).
    \item \( \hat{y}_i^{(t)} \) is the predicted probability for class \( i \).
\end{itemize}

\subsection{Simplification for One-Hot Encoded Labels}
Assume the true label \(\mathbf{y}^{(t)}\) is one-hot encoded, meaning:
\[
y_i^{(t)} = 
\begin{cases}
1 & \text{if } i = k, \\
0 & \text{otherwise},
\end{cases}
\]
where \( k \) is the index of the correct class.

Substitute \( y_i^{(t)} \) into the loss function:
\[
L^{(t)} = -\sum_{i=1}^C y_i^{(t)} \log(\hat{y}_i^{(t)})
\]
Since \( y_i^{(t)} = 0 \) for all \( i \neq k \), all terms in the sum where 
\( i \neq k \) vanish. Only the term corresponding to the correct class \( k \) 
remains:
\[
L^{(t)} = -y_k^{(t)} \log(\hat{y}_k^{(t)})
\]
But \( y_k^{(t)} = 1 \) (because \( k \) is the correct class), so:
\[
L^{(t)} = -\log(\hat{y}_k^{(t)})
\]

\subsection{Derivative of the Softmax Function}
The predicted output \(\hat{\mathbf{y}}^{(t)}\) is computed using the softmax 
function:
\[
\hat{y}_i^{(t)} = \frac{\exp(o_i^{(t)})}{\sum_{j=1}^C \exp(o_j^{(t)})}
\]
where \( o_i^{(t)} \) is the \( i \)-th element of the output vector \(\mathbf{o}^{(t)}\).

We need to compute the derivative of \(\hat{y}_k^{(t)}\) with respect to each 
element \(o_i^{(t)}\) of the output vector \(\mathbf{o}^{(t)}\).

\subsubsection{Case 1: \(i = k\) (Derivative with Respect to the Correct Class)}
When \(i = k\), we compute \(\frac{\partial \hat{y}_k^{(t)}}{\partial o_k^{(t)}}\). 
Using the quotient rule, we have:
\[
\frac{\partial \hat{y}_k^{(t)}}{\partial o_k^{(t)}} = 
\frac{\exp(o_k^{(t)}) \cdot \left( \sum_{j=1}^C \exp(o_j^{(t)}) \right) - \exp(o_k^{(t)}) \cdot \exp(o_k^{(t)})}{\left( \sum_{j=1}^C \exp(o_j^{(t)}) \right)^2}
\]
Simplify the numerator:
\[
\exp(o_k^{(t)}) \cdot \left( \sum_{j=1}^C \exp(o_j^{(t)}) - \exp(o_k^{(t)}) \right)
\]
Thus:
\[
\frac{\partial \hat{y}_k^{(t)}}{\partial o_k^{(t)}} = 
\frac{\exp(o_k^{(t)}) \cdot \left( \sum_{j=1}^C \exp(o_j^{(t)}) - \exp(o_k^{(t)}) \right)}{\left( \sum_{j=1}^C \exp(o_j^{(t)}) \right)^2}
\]
Factor out \(\exp(o_k^{(t)})\):
\[
\frac{\partial \hat{y}_k^{(t)}}{\partial o_k^{(t)}} = 
\frac{\exp(o_k^{(t)})}{\sum_{j=1}^C \exp(o_j^{(t)})} \cdot 
\frac{\sum_{j=1}^C \exp(o_j^{(t)}) - \exp(o_k^{(t)})}{\sum_{j=1}^C \exp(o_j^{(t)})}
\]
Recognize that:
\[
\frac{\exp(o_k^{(t)})}{\sum_{j=1}^C \exp(o_j^{(t)})} = \hat{y}_k^{(t)}
\]
and:
\[
\frac{\sum_{j=1}^C \exp(o_j^{(t)}) - \exp(o_k^{(t)})}{\sum_{j=1}^C \exp(o_j^{(t)})} = 1 - \hat{y}_k^{(t)}
\]
Thus:
\[
\frac{\partial \hat{y}_k^{(t)}}{\partial o_k^{(t)}} = \hat{y}_k^{(t)} (1 - \hat{y}_k^{(t)})
\]

\subsubsection{Case 2: \(i \neq k\) (Derivative with Respect to an Incorrect Class)}
When \(i \neq k\), we compute \(\frac{\partial \hat{y}_k^{(t)}}{\partial o_i^{(t)}}\). 
Using the quotient rule, we have:
\[
\frac{\partial \hat{y}_k^{(t)}}{\partial o_i^{(t)}} = 
\frac{0 \cdot \left( \sum_{j=1}^C \exp(o_j^{(t)}) \right) - \exp(o_k^{(t)}) \cdot \exp(o_i^{(t)})}{\left( \sum_{j=1}^C \exp(o_j^{(t)}) \right)^2}
\]
Simplify the numerator:
\[
-\exp(o_k^{(t)}) \cdot \exp(o_i^{(t)})
\]
Thus:
\[
\frac{\partial \hat{y}_k^{(t)}}{\partial o_i^{(t)}} = 
-\frac{\exp(o_k^{(t)}) \cdot \exp(o_i^{(t)})}{\left( \sum_{j=1}^C \exp(o_j^{(t)}) \right)^2}
\]
Factor out \(\exp(o_k^{(t)})\) and \(\exp(o_i^{(t)})\):
\[
\frac{\partial \hat{y}_k^{(t)}}{\partial o_i^{(t)}} = 
-\frac{\exp(o_k^{(t)})}{\sum_{j=1}^C \exp(o_j^{(t)})} \cdot 
\frac{\exp(o_i^{(t)})}{\sum_{j=1}^C \exp(o_j^{(t)})}
\]
Recognize that:
\[
\frac{\exp(o_k^{(t)})}{\sum_{j=1}^C \exp(o_j^{(t)})} = \hat{y}_k^{(t)}
\]
and:
\[
\frac{\exp(o_i^{(t)})}{\sum_{j=1}^C \exp(o_j^{(t)})} = \hat{y}_i^{(t)}
\]
Thus:
\[
\frac{\partial \hat{y}_k^{(t)}}{\partial o_i^{(t)}} = -\hat{y}_k^{(t)} \hat{y}_i^{(t)}
\]

\subsubsection{Summary of Softmax Derivatives}
Combining the two cases, the derivative of \(\hat{y}_k^{(t)}\) with respect to 
\(o_i^{(t)}\) is:
\[
\frac{\partial \hat{y}_k^{(t)}}{\partial o_i^{(t)}} = 
\begin{cases}
\hat{y}_k^{(t)} (1 - \hat{y}_k^{(t)}) & \text{if } i = k, \\
-\hat{y}_k^{(t)} \hat{y}_i^{(t)} & \text{if } i \neq k.
\end{cases}
\]

\subsection{Gradient of the Negative Log-Likelihood Loss}
Using the chain rule, the gradient of the loss \(L^{(t)}\) with respect to 
\(o_i^{(t)}\) is:
\[
\frac{\partial L^{(t)}}{\partial o_i^{(t)}} = 
\frac{\partial L^{(t)}}{\partial \hat{y}_k^{(t)}} 
\frac{\partial \hat{y}_k^{(t)}}{\partial o_i^{(t)}}
\]
From the negative log-likelihood loss:
\[
L^{(t)} = -\log(\hat{y}_k^{(t)})
\]
we have:
\[
\frac{\partial L^{(t)}}{\partial \hat{y}_k^{(t)}} = 
-\frac{1}{\hat{y}_k^{(t)}}
\]

Thus, the gradient is:
\[
\frac{\partial L^{(t)}}{\partial o_i^{(t)}} = 
\begin{cases}
\hat{y}_k^{(t)} - 1 & \text{if } i = k, \\
\hat{y}_i^{(t)} & \text{if } i \neq k.
\end{cases}
\]

In vector form, this can be written as:
\[
\nabla_{\mathbf{o}^{(t)}} L = \hat{\mathbf{y}}^{(t)} - \mathbf{y}^{(t)}
\]
where \(\mathbf{y}^{(t)}\) is the one-hot encoded true label.

\section{Example: \(\tau = 3\)}
We now illustrate the gradient computation with a concrete example of an 
RNN with \(\tau = 3\) time steps. We fully develop all equations for each 
time step.

\subsection{At \(t = 3\) (Final Time Step)}
At the last time step, \(\mathbf{h}^{(3)}\) only affects the current loss 
\(L^{(3)}\). Thus:
\[
\nabla_{\mathbf{h}^{(3)}} L = 
\left( \frac{\partial \mathbf{o}^{(3)}}{\partial \mathbf{h}^{(3)}} \right)^\top 
\left( \nabla_{\mathbf{o}^{(3)}} L \right)
\]
For a linear output layer:
\[
\frac{\partial \mathbf{o}^{(3)}}{\partial \mathbf{h}^{(3)}} = \mathbf{W}_{ho}
\]
Thus:
\[
\nabla_{\mathbf{h}^{(3)}} L = 
\mathbf{W}_{ho}^\top \left( \nabla_{\mathbf{o}^{(3)}} L \right)
\]
For negative log-likelihood loss:
\[
\nabla_{\mathbf{o}^{(3)}} L = \hat{\mathbf{y}}^{(3)} - \mathbf{y}^{(3)}
\]

\subsection{At \(t = 2\) (Non-Final Time Step)}
At \(t = 2\), \(\mathbf{h}^{(2)}\) affects both the current loss \(L^{(2)}\) 
and the next hidden state \(\mathbf{h}^{(3)}\). Thus:
\[
\nabla_{\mathbf{h}^{(2)}} L = 
\left( \frac{\partial \mathbf{h}^{(3)}}{\partial \mathbf{h}^{(2)}} \right)^\top 
\left( \nabla_{\mathbf{h}^{(3)}} L \right) + 
\left( \frac{\partial \mathbf{o}^{(2)}}{\partial \mathbf{h}^{(2)}} \right)^\top 
\left( \nabla_{\mathbf{o}^{(2)}} L \right)
\]
First, compute \(\left( \frac{\partial \mathbf{o}^{(2)}}{\partial \mathbf{h}^{(2)}} \right)^\top 
\left( \nabla_{\mathbf{o}^{(2)}} L \right)\):
\[
\frac{\partial \mathbf{o}^{(2)}}{\partial \mathbf{h}^{(2)}} = \mathbf{W}_{ho}
\]
Thus:
\[
\left( \frac{\partial \mathbf{o}^{(2)}}{\partial \mathbf{h}^{(2)}} \right)^\top 
\left( \nabla_{\mathbf{o}^{(2)}} L \right) = 
\mathbf{W}_{ho}^\top \left( \nabla_{\mathbf{o}^{(2)}} L \right)
\]
For negative log-likelihood loss:
\[
\nabla_{\mathbf{o}^{(2)}} L = \hat{\mathbf{y}}^{(2)} - \mathbf{y}^{(2)}
\]
Next, compute \(\left( \frac{\partial \mathbf{h}^{(3)}}{\partial \mathbf{h}^{(2)}} \right)^\top 
\left( \nabla_{\mathbf{h}^{(3)}} L \right)\):
\[
\frac{\partial \mathbf{h}^{(3)}}{\partial \mathbf{h}^{(2)}} = 
\operatorname{diag}\big(\sigma'(\mathbf{z}^{(3)})\big)\, \mathbf{W}_{hh}
\]
where \(\mathbf{z}^{(3)} = \mathbf{W}_{hh} \mathbf{h}^{(2)} + 
\mathbf{W}_{xh} \mathbf{x}^{(3)} + \mathbf{b}_h\), and \(\sigma'\) is the 
derivative of the activation function.

Thus:
\[
\nabla_{\mathbf{h}^{(2)}} L = 
\mathbf{W}_{hh}^\top \operatorname{diag}\big(\sigma'(\mathbf{z}^{(3)})\big) 
\left( \nabla_{\mathbf{h}^{(3)}} L \right) + 
\mathbf{W}_{ho}^\top \left( \nabla_{\mathbf{o}^{(2)}} L \right)
\]

\subsection{At \(t = 1\) (Non-Final Time Step)}
At \(t = 1\), \(\mathbf{h}^{(1)}\) affects both the current loss \(L^{(1)}\) 
and the next hidden state \(\mathbf{h}^{(2)}\). Thus:
\[
\nabla_{\mathbf{h}^{(1)}} L = 
\left( \frac{\partial \mathbf{h}^{(2)}}{\partial \mathbf{h}^{(1)}} \right)^\top 
\left( \nabla_{\mathbf{h}^{(2)}} L \right) + 
\left( \frac{\partial \mathbf{o}^{(1)}}{\partial \mathbf{h}^{(1)}} \right)^\top 
\left( \nabla_{\mathbf{o}^{(1)}} L \right)
\]
First, compute \(\left( \frac{\partial \mathbf{o}^{(1)}}{\partial \mathbf{h}^{(1)}} \right)^\top 
\left( \nabla_{\mathbf{o}^{(1)}} L \right)\):
\[
\frac{\partial \mathbf{o}^{(1)}}{\partial \mathbf{h}^{(1)}} = \mathbf{W}_{ho}
\]
Thus:
\[
\left( \frac{\partial \mathbf{o}^{(1)}}{\partial \mathbf{h}^{(1)}} \right)^\top 
\left( \nabla_{\mathbf{o}^{(1)}} L \right) = 
\mathbf{W}_{ho}^\top \left( \nabla_{\mathbf{o}^{(1)}} L \right)
\]
For negative log-likelihood loss:
\[
\nabla_{\mathbf{o}^{(1)}} L = \hat{\mathbf{y}}^{(1)} - \mathbf{y}^{(1)}
\]
Next, compute \(\left( \frac{\partial \mathbf{h}^{(2)}}{\partial \mathbf{h}^{(1)}} \right)^\top 
\left( \nabla_{\mathbf{h}^{(2)}} L \right)\):
\[
\frac{\partial \mathbf{h}^{(2)}}{\partial \mathbf{h}^{(1)}} = 
\operatorname{diag}\big(\sigma'(\mathbf{z}^{(2)})\big)\, \mathbf{W}_{hh}
\]
where \(\mathbf{z}^{(2)} = \mathbf{W}_{hh} \mathbf{h}^{(1)} + 
\mathbf{W}_{xh} \mathbf{x}^{(2)} + \mathbf{b}_h\), and \(\sigma'\) is the 
derivative of the activation function.

Thus:
\[
\nabla_{\mathbf{h}^{(1)}} L = 
\mathbf{W}_{hh}^\top \operatorname{diag}\big(\sigma'(\mathbf{z}^{(2)})\big) 
\left( \nabla_{\mathbf{h}^{(2)}} L \right) + 
\mathbf{W}_{ho}^\top \left( \nabla_{\mathbf{o}^{(1)}} L \right)
\]

\section{Conclusion}
The gradient computation in RNNs involves both direct and indirect 
effects of hidden states on the loss. The recursive nature of the 
gradient equation ensures that all paths of influence are accounted for. 
This document provides a detailed explanation of the general case, 
splitting the computations for the final step and non-final steps, and 
fully develops all equations for a concrete example with \(\tau = 3\). 
It also includes a detailed derivation of the Negative Log-Likelihood 
Loss and its gradient, including all computational steps for the derivative 
of the softmax function.
\end{notechapter}
